\section{TRABAJOS RELACIONADOS}

En el primer trabajo de Kalabarige, la metodología parte de una selección híbrida de características, pues se entrenan tres algoritmos de boosting (XGBoost, CatBoost y LightGBM), se extrae la importancia que cada uno le asigna a las variables, se promedia esa importancia entre los tres modelos y se conservan las K variables mejor puntuadas, con K variando entre el 66\% y el 86\% del total según el dataset. Sobre ese subconjunto de características se entrena un modelo de apilamiento de tres capas construido con algoritmos de boosting, donde la primera capa usa XGBoost, LightGBM, CatBoost y AdaBoost, la segunda descarta AdaBoost y conserva los tres restantes, y la tercera usa XGBoost como meta aprendiz final. El proceso se repite en cuatro fases, con todas las características sin balanceo, con las K seleccionadas, sin balanceo y ambas variantes con balanceo mediante sobremuestreo aleatorio para medir el efecto de cada decisión por separado \cite{kalabarige2023boosting}. \newline

En el trabajo previo de los mismos autores, la metodología es más sencilla en cuanto a la etapa de features, ya que no hay selección, sino que se usan todas las variables del dataset. El aporte metodológico está en la arquitectura de apilamiento, pues la primera capa combina cinco clasificadores heterogéneos (MLP, KNN, Random Forest, Regresión Logística y XGBoost), de esos cinco se seleccionan los tres de mejor desempeño para conformar la segunda capa, y la tercera usa nuevamente XGBoost como meta aprendiz. Antes de entrenar, los datasets con clases desbalanceadas se balancean mediante sobremuestreo aleatorio (ROS), duplicando instancias de la clase minoritaria hasta igualarla con la mayoritaria, y el modelo se evalúa comparando su desempeño con y sin ese balanceo \cite{kalabarige2022multilayer}. \newline

Duy aborda la metodología desde un ángulo distinto, centrado en poner a prueba la robustez de los modelos más que en optimizar su exactitud. Por lo que construyen un framework llamado AWG que combina una red generativa adversaria tipo Wasserstein con un algoritmo de generación de dominios maliciosos (Charbot), de forma que un generador aprende a producir URLs adversarias que engañan a un discriminador entrenado para imitar el comportamiento de un detector de caja negra. Con las muestras generadas atacan quince modelos distintos, de tipo ML, ensamble, aprendizaje profundo y multimodal, bajo un escenario de ataque de transferencia, donde el atacante no conoce la arquitectura interna del modelo víctima. Como segunda etapa metodológica, proponen una estrategia de defensa por reentrenamiento, donde seleccionan las muestras adversarias de mayor calidad (las que lograron evadir un clasificador de dominios) y las mezclan con el dataset original para reentrenar los modelos víctima, midiendo la mejora en su tasa de detección \cite{duy2024study}. \newline

Ahmad no propone un modelo propio, sino que su metodología es la de una revisión sistemática de literatura. Parten de una búsqueda con la palabra clave "Web Phishing Detection", que arrojó cerca de 3,500 artículos, de esos artículos van filtrando por especificidad y por año de publicación (2020-2024) hasta quedarse con 150 artículos, divididos en 100 de investigación original, 30 encuestas y 20 revisiones. Sobre ese conjunto organizan la información en categorías, modos de ataque, técnicas de detección basadas en listas, en heurísticas, en el lado del cliente y en modelos de IA, y comparan el desempeño reportado por cada estudio usando la exactitud como métrica común, lo que les permite construir un mapa comparativo de qué modelos y qué técnicas de extracción de características han dado mejores resultados en la literatura reciente \cite{ahmad2025across}. \newline

Por último, Alsariera fija un único algoritmo base (Extra-Trees) y varían la estrategia de meta aprendizaje sobre él, dando lugar a cuatro modelos, AdaBoost-Extra Tree, Bagging-Extra Tree, Rotation Forest-Extra Tree y LogitBoost-Extra Tree. Los cuatro se entrenan y validan mediante validación cruzada de 10 particiones sobre el dataset de phishing de UCI, particionando los datos en diez grupos, entrenando sobre nueve y probando sobre el restante en cada iteración. El desempeño de cada modelo se mide con exactitud, tasa de falsos positivos, tasa de falsos negativos y F-measure, y luego se compara contra quince métodos previos publicados entre 2014 y 2020 para ubicar los resultados dentro del estado del arte \cite{alsariera2020ai}. \newline

La metodología de este proyecto se apoya en el enfoque de Kalabarige \cite{kalabarige2023boosting} de usar varios modelos de boosting para promediar la importancia de las variables antes de entrenar el modelo final, pero la simplifica a una sola capa de apilamiento, con XGBoost y Random Forest como clasificadores base y un segundo XGBoost como meta clasificador. El detalle completo del pipeline, del proceso de entrenamiento y de la estrategia de evaluación se desarrolla en la siguiente sección.